% 本文件由 03_代码/p3_tables.py 自动生成，请勿手工修改。
% 数据源：06_支撑材料/p3_analysis.json

\begin{table}[htbp]
  \centering
  \caption{八种预报组合在 2025 年 2 月 1 日--12 月 31 日的实际费用（参数固定为 $\alpha=0.6,\rho=0.5,\lambda=1.0,W=28$，评价区间、负载模型、执行规则与计费口径完全一致）}
  \label{tab:p3-combo}
  \small
  \begin{tabular}{lrrrrr}
    \toprule
    使用的预报时刻 & 计划调整费/元 & 紧急费/元 & 合计费用/元 & 相对基准 & 期末储电量/kWh \\
    \midrule
    仅 0{:}00 & 12,595,136 & 2,382,673 & \textbf{14,977,809} & --- & 6,007.5 \\
    6{:}00 & 12,767,094 & 1,712,866 & \textbf{14,479,959} & -3.32\% & 6,007.5 \\
    12{:}00 & 12,767,834 & 1,264,057 & \textbf{14,031,892} & -6.32\% & 6,007.5 \\
    6{:}00+12{:}00 & 12,874,840 & 1,006,040 & \textbf{13,880,879} & -7.32\% & 6,007.5 \\
    18{:}00 & 13,187,014 & 612,575 & \textbf{13,799,589} & -7.87\% & 5,798.2 \\
    6{:}00+18{:}00 & 13,259,141 & 273,874 & \textbf{13,533,015} & -9.65\% & 5,798.2 \\
    12{:}00+18{:}00 & 13,033,851 & 546,355 & \textbf{13,580,206} & -9.33\% & 5,798.2 \\
    6{:}00+12{:}00+18{:}00 & 13,145,690 & 272,936 & \textbf{13,418,626} & -10.41\% & 5,798.2 \\
    \midrule
    \textbf{主策略（全用）} & 13,145,690 & 272,936 & \textbf{13,418,626} & -10.41\% & 5,798.2 \\
    \bottomrule
  \end{tabular}
  \par\vspace{2pt}
  \begin{minipage}{\textwidth}\footnotesize
    注：所有方案都使用 0:00 发布的预报；组合列出的时刻是\textbf{额外}使用的版本。首行“仅 0:00”是关闭全部调整的对照，其购电量不随日内预报变化。“相对基准”以该对照为分母，负值表示比不做调整更省。合计费用已含计划部分，不能再与计划调整费、紧急费相加。
  \end{minipage}
\end{table}

\begin{table}[htbp]
  \centering
  \caption{附件 3 四个发布版本在 2025 年 2 月 1 日--12 月 31 日的光伏出力预报精度（单位：kW；把整点预报因果内插到十分钟后与实际光伏逐段比较）}
  \label{tab:p3-skill}
  \small
  \begin{tabular}{llrrr}
    \toprule
    预报版本 & 当天可用时段 & 可用时段数 & 覆盖段绝对误差 & 其中 $12{:}00$--$18{:}00$ \\
    \midrule
    0{:}00 发布 & 0{:}00--24{:}00 & 48,096 & 206.6 & 445.8 \\
    6{:}00 发布 & 6{:}00--24{:}00 & 36,072 & 180.0 & 304.9 \\
    12{:}00 发布 & 12{:}00--24{:}00 & 24,048 & 98.6 & 181.3 \\
    18{:}00 发布 & 18{:}00--24{:}00 & 12,024 & 15.9 & --- \\
    \bottomrule
  \end{tabular}
  \par\vspace{2pt}
  \begin{minipage}{\textwidth}\footnotesize
    注：四个版本的覆盖时段长度不同，因此采用绝对误差。$18{:}00$ 版仅覆盖光伏接近零的夜间，相对误差会受较小分母影响。在四版共同覆盖的 $18{:}00$--$24{:}00$ 时段，误差均为 $15.7$--$15.9$~kW；$18{:}00$ 版不覆盖白天，相应单元格留空。预报版本的经济作用还取决于其对应的调整时段。
  \end{minipage}
\end{table}

\begin{table}[htbp]
  \centering
  \caption{主策略参数的单变量扰动（每次只动一个参数，其余取 $\alpha=0.6,\rho=0.5,\lambda=1.0,W=28$，评价区间同前）}
  \label{tab:p3-param}
  \small
  \begin{tabular}{lrrrr}
    \toprule
    参数取值 & 计划调整费/元 & 紧急费/元 & 合计费用/元 & 相对主策略 \\
    \midrule
    主策略 $\alpha=0.6,\rho=0.5,\lambda=1.0,W=28$ & 13,145,690 & 272,936 & \textbf{13,418,626} & --- \\
    \midrule
    $W=$14 & 13,165,617 & 297,334 & 13,462,951 & +0.33\% \\
    $W=$21 & 13,158,472 & 282,713 & 13,441,185 & +0.17\% \\
    $\alpha=$0.5,\;$\rho=$0.35 & 13,053,920 & 511,608 & 13,565,528 & +1.09\% \\
    $\alpha=$0.5,\;$\rho=$0.5 & 13,052,126 & 516,895 & 13,569,021 & +1.12\% \\
    $\alpha=$0.5,\;$\rho=$0.7 & 13,040,807 & 545,535 & 13,586,342 & +1.25\% \\
    $\alpha=$0.6,\;$\rho=$0.35 & 13,146,530 & 271,409 & 13,417,938 & -0.01\% \\
    $\alpha=$0.6,\;$\rho=$0.7 & 13,139,824 & 286,336 & 13,426,160 & +0.06\% \\
    $\alpha=$0.7,\;$\rho=$0.35 & 13,327,609 & 137,247 & 13,464,856 & +0.34\% \\
    $\alpha=$0.7,\;$\rho=$0.5 & 13,327,039 & 137,873 & 13,464,912 & +0.34\% \\
    $\alpha=$0.7,\;$\rho=$0.7 & 13,325,109 & 141,702 & 13,466,811 & +0.36\% \\
    $\alpha=$0.8,\;$\rho=$0.35 & 13,648,143 & 67,335 & 13,715,479 & +2.21\% \\
    $\alpha=$0.8,\;$\rho=$0.5 & 13,647,810 & 67,850 & 13,715,659 & +2.21\% \\
    $\alpha=$0.8,\;$\rho=$0.7 & 13,646,892 & 69,426 & 13,716,318 & +2.22\% \\
    $\lambda=$0.0 & 13,134,961 & 300,666 & 13,435,627 & +0.13\% \\
    $\lambda=$0.5 & 13,135,859 & 272,916 & 13,408,774 & -0.07\% \\
    \bottomrule
  \end{tabular}
  \par\vspace{2pt}
  \begin{minipage}{\textwidth}\footnotesize
    注：$\alpha$ 为净负荷预测的风险分位数，$\rho$ 为实时放电储备线比例，$\lambda$ 为日末储电量回归目标的引导权重，$W$ 为分位数所用的历史窗口天数。表中省略了与主策略完全相同的取值行。
  \end{minipage}
\end{table}

\begin{table}[htbp]
  \centering
  \caption{计费口径与风险分位数口径的敏感性（模型、参数、评价区间均与主策略相同，只换这一处口径）}
  \label{tab:p3-billing}
  \small
  \begin{tabular}{lrrrr}
    \toprule
    口径 & 计划调整费/元 & 紧急费/元 & 合计费用/元 & 相对主策略 \\
    \midrule
    主策略：退款口径 $+$ 跨时段合并分位数 & 13,145,690 & 272,936 & \textbf{13,418,626} & --- \\
    调减不退款口径 & 13,357,217 & 261,727 & 13,618,944 & +1.49\% \\
    逐时段分位数 & 13,082,057 & 302,772 & 13,384,829 & -0.25\% \\
    \bottomrule
  \end{tabular}
  \par\vspace{2pt}
  \begin{minipage}{\textwidth}\footnotesize
    注：题目未明确调减部分的原价是否退回。主策略采用退款口径，即调减部分按 $0.5p_t$ 结算；不退款口径仍支付原计划费用。本表并列报告两种计费结果。第二项比较跨时段合并分位数与逐时段分位数，用于衡量风险估计粒度的影响。
  \end{minipage}
\end{table}

\begin{table}[htbp]
  \centering
  \caption{参数的 1 月标定记录：标定面，以及用 1 月最优参数直接外推的检验}
  \label{tab:p3-calib}
  \small
  \begin{tabular}{lrrrr}
    \toprule
    风险分位数 $\alpha$ & $\rho=0.35$ & $\rho=0.5$ & $\rho=0.7$ & 该行最小 \\
    \midrule
    0.4 & \textbf{1,802,432} & 1,811,016 & 1,822,471 & 1,802,432 \\
    0.5 & \textbf{1,770,321} & 1,778,833 & 1,790,196 & 1,770,321 \\
    0.6 & \textbf{1,764,349} & 1,772,822 & 1,784,039 & 1,764,349 \\
    0.7 & \textbf{1,782,106} & 1,790,579 & 1,801,805 & 1,782,106 \\
    0.8 & \textbf{1,830,510} & 1,838,983 & 1,850,209 & 1,830,510 \\
    \midrule
    窗口 $W=14$ & \multicolumn{3}{c}{1,771,570} &  \\
    窗口 $W=21$ & \multicolumn{3}{c}{1,772,674} &  \\
    窗口 $W=28$ & \multicolumn{3}{c}{1,772,822} & （采用） \\
    \midrule
    $\lambda=0.0$ & \multicolumn{3}{c}{1,774,575} & （日末储电仅 1,806~kWh） \\
    $\lambda=0.5$ & \multicolumn{3}{c}{1,769,148} &  \\
    $\lambda=1.0$ & \multicolumn{3}{c}{1,772,822} & （采用） \\
    \bottomrule
  \end{tabular}
  \par\vspace{2pt}
  \begin{minipage}{\textwidth}\footnotesize
    上半部分为 1 月标定面（$W=28$，单位：元），下半部分为 $W$ 与 $\lambda$ 的单变量扫描。标定面在 $\alpha$ 方向有明确的极小点 $\alpha=0.6$，在 $\rho$ 与 $W$ 方向则相当平坦：$\rho\in[0.35,0.5]$、$W\in[14,28]$ 的全部取值都落在同一水平内，彼此相差不到 $1\%$。
    \textbf{样本外检验：}1 月标定面上的最小值为 $(\alpha,\rho,W)=(0.6,0.35,14)$，1 月费用 $1,763,097$~元；本文采用的 $(0.6,0.5,28)$ 对应 $1,772,822$~元，相差 $9,725$~元。在完整评价区间内，两组参数的费用分别为 $13,461,840$~元和 $13,418,626$~元，后者低 $43,213$~元。考虑标定面在 $\rho$、$W$ 方向较平坦，采用中心取值 $(0.6,0.5,28)$。$\lambda=1.0$ 与 $\lambda=0.5$ 在评价区间内相差 $0.07\%$（$9{,}852$~元）；$\lambda=0$ 时日末储电量为 $1{,}806$~kWh，因此保留非零日末储备权重。
  \end{minipage}
\end{table}

\begin{table}[htbp]
  \centering
  \caption{九项检查在主策略全部 334 个交付日上的实测结果}
  \label{tab:p3-check}
  \small
  \begin{tabular}{llr}
    \toprule
    检查项 & 实测最坏值 & 未通过天数 \\
    \midrule
    逐段能量平衡 & $2.27e-13$ kWh & 0 \\
    储能状态递推 & $0.00e+00$ kWh & 0 \\
    计划调整费复算 & $0.00e+00$ 元 & 0 \\
    紧急购电费复算 & $0.00e+00$ 元 & 0 \\
    储电量上下界 & 结构上恒成立 & 0 \\
    充放电功率上限 & 结构上恒成立 & 0 \\
    同时充放电 & 结构上恒成立 & 0 \\
    购电量非负 & 结构上恒成立 & 0 \\
    跨日状态衔接 & 结构上恒成立 & 0 \\
    \midrule
    合计 & --- & \textbf{0} \\
    \bottomrule
  \end{tabular}
  \par\vspace{2pt}
  \begin{minipage}{\textwidth}\footnotesize
    注：九项检查在全部交付日上一次通过，无任何残差、越界或复算不符。
    前四项由独立复算得到最坏值，后五项由线性规划的约束边界与实时执行规则结构性地保证，故逐日实测值恒为零。
  \end{minipage}
\end{table}
